\section*{Bab \ref{ch:b2:wilson-cycle} --- Masif Tua dan Siklus Wilson}

\begin{solution}{exo:b2:wilson-cycle:1}
(a) $600 - 400 = \qty{200}{Ma}$, dari pembentukan celah di tepi-tepinya sampai tumbukan terakhir.
(b) $480 - 400 = \qty{80}{Ma}$. (c) Orogenesis Kaledonia adalah serangkaian fase yang tersebar sepanjang
sekitar 200 juta tahun, dan setiap fase dimulai pada waktu yang berbeda di tempat yang berbeda: kata itu
menunjuk sebuah samudra dan sabuknya, bukan suatu saat. Umur diberikan oleh skala stratigrafi atau dalam
juta tahun.
\end{solution}

\begin{solution}{exo:b2:wilson-cycle:2}
(c) celah Afrika Timur, pembentukan celah; (a) Laut Merah, samudra muda; (d) Atlantik, samudra luas
dengan tepi pasif; (e) tepi Pasifik Amerika Selatan, subduksi dan
penutupan samudra; (b) Himalaya, tumbukan; (f) Bretagne, sabuk yang telah terkikis rata, yang setelahnya
pembentukan celah dapat dimulai lagi.
\end{solution}

\begin{solution}{exo:b2:wilson-cycle:3}
(a) Setelah pengendapan lapisan Silur awal, jadi setelah \qty{443.1}{Ma}, dan sebelum
pengendapan lapisan Devon Akhir, jadi sebelum \qty{358.9}{Ma}, akhir Devon.
(b) Kurang dari $443.1 - 358.9 = \qty{84.2}{Ma}$; sebenarnya lebih singkat lagi, karena lapisan bawah lebih muda
daripada dasar Silur dan lapisan atas lebih tua daripada puncak Devon. (c) Pengendapan
turbidit kelabu di samudra; pelipatan dan belahan dalam tumbukan-tumbukan Kaledonia; pengangkatan dan
erosi menjadi permukaan daratan yang tidak rata; pengendapan breksi dan konglomerat merah; pemiringan
landai dalam orogenesis Variscan; erosi masa kini, yang memotong tebing itu.
\end{solution}

\begin{solution}{exo:b2:wilson-cycle:4}
(a) Keduanya: sebuah blok kerak benua yang terpisah dari benua-benua oleh samudra selagi
menyeberangi Iapetus, dan, setelah teraduk, sebuah blok berbatas sesar yang sejarahnya berbeda dari sejarah
blok-blok tetangganya. (b) Bukan keduanya: sebuah benua, bukan blok yang ditambahkan ke benua. (c) \omterm{def:b2:wilson-cycle:terrane}{Terrane}, bukan
\omterm{def:b2:wilson-cycle:terrane}{mikrobenua}: keraknya adalah kerak busur, bukan pecahan benua. (d) Bukan keduanya: sutura itu adalah
batas antara \omterm{def:b2:wilson-cycle:terrane}{terrane-terrane}, jejak samudra yang telah menutup.
\end{solution}

\begin{solution}{exo:b2:wilson-cycle:5}
(a) $3000 - 2600 = 400$, $2600 - 1500 = 1100$, $1500 - 875 = 625$, $875 - 600 = 275$, $600 -
262.5 = 337.5$; rata-rata $2737.5/5 = \qty{547.5}{Ma}$. (b) $1100/275 = 4.0$. (c) $(1500 -
262.5)/3 = \qty{412.5}{Ma}$. Rata-ratanya bergeser seperempat bila titik-titik yang diragukan dibuang,
dan jarak-jaraknya bervariasi empat kali lipat: suatu perulangan kasar, bukan periode.
\end{solution}

\begin{solution}{exo:b2:wilson-cycle:6}
(a) $3000/(70 - 25) = \qty{66.7}{Ma}$. (b) $3000/70 = \qty{42.9}{Ma}$. (c) Setelah punggungan
tersubduksi, dasar samudra yang tersisa adalah bagian lempeng jauh, yang membawa benua
jauh: laju penutupan adalah kecepatan lempeng itu relatif terhadap parit.
\end{solution}

\begin{solution}{exo:b2:wilson-cycle:7}
(a) $700/8 = \qty{87.5}{km}$; $700/30 = \qty{23.3}{km}$. (b) Sabuk Variscan: pada
\qty{30}{\celsius} per kilometer kerak tengah meleleh; pada \qty{8}{\celsius} per kilometer
suhu pelelehan terletak di bawah dasar kerak, bahkan kerak yang menebal. (c) Gradien rendah itu
menggambarkan batuan yang terbawa turun dengan cepat di zona subduksi, sebelum sempat memanas. Batuan yang
tetap berada di kedalaman dalam kerak yang menebal menghangat menuju geoterm normal selama jutaan tahun, sehingga dapat
mencapai pelelehan pada kedalaman yang lebih kecil.
\end{solution}

\begin{solution}{exo:b2:wilson-cycle:8}
(a) $\tau\ln(3/0.3) = 26.3\times2.303 = \qty{60.6}{Ma}$. (b) $290 - 60.6 = \qty{229}{Ma}$, pada
Trias Akhir. (c) $\tau = (290 - 251.9)/\ln10 = \qty{16.5}{Ma}$. (d) Keruntuhan ekstensional
kerak yang menebal, yang menipiskannya dan menurunkan permukaannya secara langsung; juga laju
erosi yang lebih tinggi, iklim yang lebih basah atau batuan yang lebih lunak, atau elevasi awal yang lebih rendah.
\end{solution}

\begin{solution}{exo:b2:wilson-cycle:9}
(a) Sutura, jejak Iapetus, terletak di daratan Inggris dan Irlandia, di timur
Atlantik; kerak Laurentia, dari sisi seberang Iapetus, terletak di dekat Rockall di sisi
Eropa. Celah baru itu memotong di barat sutura lama: ia membuka di dekatnya, bukan di sepanjangnya.
(b) Di utara sutura, batuan dasar dan penutup tepi Laurentia; di selatannya, batuan dasar dan penutup
Avalonia, dengan sejarah yang berbeda dan, pada Paleozoikum awal, fauna yang berbeda. (c) Pembentukan celah
datang setelah 80 juta tahun, lebih dari 30 sampai 50 juta tahun yang setelahnya zona yang telah mendingin mungkin telah menjadi
kuat; namun Atlantik Tengah memang mengikuti Appalachia. Panas suatu sabuk memudar; sesar-sesarnya,
fabriknya, dan komposisinya tetap ada dan masih dapat memandu sebuah celah.
\end{solution}

\begin{solution}{exo:b2:wilson-cycle:10}
(a) $70 - 40/2 = \qty{50}{mm}$ per tahun. (b) $\qty{1500}{km}/\qty{50}{km.Ma^{-1}} =
\qty{30}{Ma}$. (c) Samudra menyempit dengan laju $70 - 40 = \qty{30}{mm}$ per tahun: $3000 - 30\times30 =
\qty{2100}{km}$. (d) Litosfer muda yang panas lalu tiba di parit, dan tidak ada lagi slab tua yang dingin
yang melepaskan air di bawah busur: magmatisme busur berubah atau berhenti.
\end{solution}

\begin{solution}{exo:b2:wilson-cycle:11}
(a) $h = \Delta T(\rho_m - \rho_c)/\rho_m = 25\times600/3300 = \qty{4.5}{km}$. (b) Seluruh tebal tambahan
\qty{25}{km}, 5.5 kali tinggi jajaran pegunungan itu. (c) $2700\times9.8\times25\,000 =
\qty{6.6e8}{Pa}$, sekitar \qty{0.66}{GPa}. (d) Pada sekitar \qty{0.7}{GPa} dan suhu
akar kerak yang menebal, beberapa ratus derajat, fasies amfibolit atau fasies
granulit: gneis dan \omterm{def:b2:crustal-melting:migmatite}{migmatit} masif-masif tua.
\end{solution}

\begin{solution}{exo:b2:wilson-cycle:12}
(a) $S = 30\times10 = \qty{300}{km}$; $W_0 = 300 + 300 = \qty{600}{km}$; $T = 35\times600/300 =
\qty{70}{km}$. (b) $h = 35\times600/3300 = \qty{6.36}{km}$; $r = 4.5h = \qty{28.6}{km}$.
(c) $\Delta F = 2700\times9.8\times6364\times(35\,000 + 3182 + 14\,318) =
\qty{8.8e12}{N.m^{-1}}$, 1.4 kali energi Tibet. (d) Sabuk seperti itu lebih berat daripada Tibet, yang sudah
menyebar: \omterm{def:b2:himalaya-andes:collapse}{keruntuhan gravitasi}, melalui ekstensi di wilayah-wilayah tingginya dan sesar naik di tepi-tepinya.
\end{solution}

\begin{solution}{pb:b2:wilson-cycle:1}
\textbf{1.} $720/(540 - 500) = \qty{18}{km}$ per juta tahun, \qty{18}{mm} per tahun.
\textbf{2.} \qty{36}{mm} per tahun: punggungan berpemekaran lambat (20 sampai \qty{55}{mm} per tahun).
\textbf{3.} $36\times(580 - 540) = \qty{1440}{km}$.
\textbf{4.} $36\times(580 - 500) = \qty{2880}{km}$.
\textbf{5.} Pemekaran simetris dengan laju tetap dari lebar nol pada \qty{580}{Ma}, dan tidak ada
subduksi sebelum \qty{500}{Ma}.
\textbf{6.} \qty{80}{Ma}, di kedua tepi samudra, bersebelahan dengan tepi-tepi benua. Itulah litosfer
samudra yang paling dingin, paling tebal, dan paling padat di samudra itu, yang paling siap tenggelam.
\textbf{7.} $64 - 36 = \qty{28}{mm}$ per tahun.
\textbf{8.} $2880/28 = \qty{102.9}{Ma}$ setelah \qty{500}{Ma}: pada \qty{397.1}{Ma}.
\textbf{9.} Punggungan, yang berjarak \qty{1440}{km} dari parit, mendekatinya dengan laju $64 - 18 =
\qty{46}{mm}$ per tahun: setelah $1440/46 = \qty{31.3}{Ma}$, pada \qty{468.7}{Ma}.
\textbf{10.} $2880 - 28\times31.3 = \qty{2004}{km}$.
\textbf{11.} $2004/40 = \qty{50.1}{Ma}$ kemudian: pada $468.7 - 50.1 = \qty{418.6}{Ma}$.
\textbf{12.} Di dalam rentang 425 sampai \qty{400}{Ma} tumbukan Skandia.
\textbf{13.} $40\times8 = \qty{320}{km}$.
\textbf{14.} $W_0 = 400 + 320 = \qty{720}{km}$; $T = 35\times720/400 = \qty{63.0}{km}$.
\textbf{15.} $h = 28\times600/3300 = \qty{5.09}{km}$; $r = 4.5\times5.09 = \qty{22.9}{km}$.
\textbf{16.} $\Delta F = 2700\times9.8\times5091\times(35\,000 + 2545 + 11\,455) =
\qty{6.60e12}{N.m^{-1}}$, sekitar \qty{8}{\percent} lebih besar daripada energi Tibet.
\textbf{17.} Berhenti naik dan menyebar: ekstensi dan sesar normal di bagian dalamnya yang tinggi,
sesar naik di tepi-tepinya, seperti yang dilakukan Tibet kini.
\textbf{18.} $f = 600/3300 = 0.182$; $k = 2.303/11 = 0.209$ per juta tahun; $\tau = 1/(kf)
= \qty{26.3}{Ma}$.
\textbf{19.} $3.0\,e^{-30/26.3} = \qty{0.96}{km}$.
\textbf{20.} $26.3\ln(3.0/0.5) = \qty{47.1}{Ma}$ setelah \qty{410}{Ma}: pada \qty{362.9}{Ma}.
\textbf{21.} Kerak tambahan adalah $h/f$: $3.0/0.182 = \qty{16.5}{km}$ turun menjadi $0.5/0.182 =
\qty{2.75}{km}$, sehingga \qty{13.75}{km} batuan tersingkir.
\textbf{22.} $2700\times9.8\times13\,750 = \qty{3.6e8}{Pa}$, sekitar \qty{0.36}{GPa}.
\textbf{23.} Ya: pada Devon Akhir elevasi rata-rata sabuk model itu tinggal sekitar \qty{1}{km}
dan masih turun, dengan relief yang tersisa di antara lembah dan punggung bukitnya; kerikil merah dapat
terkumpul di bagian-bagian rendahnya jauh sebelum \qty{358.9}{Ma}, seperti di Siccar Point, tempat kerikil itu bertumpu di atas permukaan
yang berelief.
\textbf{24.} \textbf{\qty[text-series-to-math=true]{2880}{km}}, pada \qty{500}{Ma}, ketika
subduksi dimulai.
\end{solution}
